\section{Experimental Setup}
\label{sec:setup}

\subsection{Document Collection and Queries}
\label{subsec:collection}

\paragraph{Document corpus.}
The document collection consists of 10{,}000 scientific publications 
provided by the CLEF~2026 CheckThat!~Lab Task~1 organisers.
Each publication is described by the following metadata fields: 
\texttt{title}, \texttt{abstract}, \texttt{authors}, \texttt{venue}, 
and a unique identifier \texttt{pubkey}.
In our system, we index only the \texttt{title} and \texttt{abstract} 
fields, as these carry the core semantic content of a paper and are 
available for all documents in the corpus.

\paragraph{Query sets.}
The queries consist of short social media posts (tweets), partitioned 
into three official splits: training, development, and test.
Each query is provided as a \texttt{text} field containing the raw 
social media content, along with a \texttt{pubkey} field that identifies 
the single relevant document in the corpus.
This yields a \textit{one-to-one relevance mapping}: each query is 
associated with exactly one ground-truth document, making this a 
highly precise but recall-agnostic retrieval setting.

Table~\ref{tab:collection} summarises the statistics of the collection 
and query splits.
\begin{table}[h]
  \caption{Collection and query statistics.}
  \label{tab:collection}
  \centering
  \begin{tabular}{lrrr}
    \toprule
    Split & English & German & French \\
    \midrule
    Train queries & 14977 & 1460 & 2807 \\
    Dev queries   & 3905 & 386 & 702 \\
    Test queries  & 6080 & 876 & 1220 \\
    \midrule
    Corpus documents & \multicolumn{3}{c}{10{,}000} \\
    \bottomrule
  \end{tabular}
\end{table}

\subsection{Evaluation Measures}
\label{subsec:measures}

We evaluate retrieval quality using the official metric \ac{MRR}@5, which rewards
systems that place the single relevant document in the top-5 ranked positions.
In addition, we report \ac{MAP}, \ac{NDCG}@10, Precision@1, and Recall@1000 to
characterise different aspects of ranking quality.
Evaluation is performed with the \texttt{ranx} library using relevance
judgements (\texttt{qrels}) derived from the ground-truth \texttt{pubkey} labels
provided with each query.

\subsection{Repository}
\label{subsec:repo}

The source code is publicly available at:
\url{https://bitbucket.org/upd-dei-stud-prj/seupd2526-dspy/src/master/}.
%
The repository is organised as follows:
\begin{itemize}
  \item \texttt{pylucene/src/} — entry point orchestrator (\texttt{runner.py}) and data acquisition scripts.
  \item \texttt{pylucene/src/Index/} — core PyLucene and FAISS indexing pipelines (\texttt{index.py}, \\ \texttt{index\_faiss.py}).
  \item \texttt{pylucene/src/Search/} — core execution scripts for the four retrieval methodologies \\(\texttt{search\_v1.py} to \texttt{search\_v4.py}).
  \item \texttt{pylucene/src/Parse/} — data pre-processing and formatting classes for documents and queries.
  \item \texttt{pylucene/src/Evaluation/} \& \texttt{utils/} — scoring (\texttt{evaluate.py}), report generation, ground truth creation, and submission formatting.
  \item \texttt{experiment/} — generated artefacts (index, runs, qrels); not committed to version control.
\end{itemize}

\subsection{Hardware and Software}
\label{subsec:hardware}

Experiments were conducted on three machines. All ran the retrieval pipeline
inside a Docker container built from the official \texttt{coady/pylucene} image,
which bundles PyLucene~9 (Lucene~9.x) on Python~3.11 (to check).

\begin{table}[h]
  \caption{Hardware specifications of the development machines.}
  \label{tab:hardware}
  \centering
  \begin{tabular}{llll}
    \toprule
    Component & Machine~1 & Machine~2 & Machine~3 \\
    \midrule
    CPU & Intel Core i5-13420H  & AMD Ryzen AI 7 350    & Intel Core i7-6500U   \\
        & (13th Gen)            &                       & @ 2.50\,GHz           \\
    RAM & 32\,GB DDR5 5600\,MHz & 16\,GB DDR5 5600\,MHz & 12\,GB DDR4 2133\,MHz \\
    GPU & Intel UHD Graphics    & AMD Radeon 860M +     & NVIDIA GeForce 920MX  \\
        &                       & NVIDIA RTX~5060 Laptop&                       \\
    OS  & Windows~11 Home       & Windows~11 Home       & Windows~10 Home       \\
    \bottomrule
  \end{tabular}
\end{table}

The Docker environment also exposes a Jupyter notebook server on port~8888 for Exploratory Data Analysis. The GPU of Machine~2 was used to accelerate the neural cross-encoder re-ranker (\texttt{rerank.py}) and FAISS dense retrieval experiments (\texttt{index\_faiss.py} and \texttt{search\_faiss.py}). All BM25 indexing , and \ac{PRF}  are fully reproducible on \ac{CPU} and were run on all three machines.
